%% ma: L12-B08-0192
%% lop: 12
%% chuong: 2
%% bai: 8
%% dang: 12.8.37
%% loai: TN4
%% muc: TH
%% dap_an: "A"
%% nguon: "CĐ2.6.pdf giáo viên gửi 27/10/2026 (M07, Lê Minh Khôi – Nguyễn Anh Khôi, Chương II Chủ đề 6 Ứng dụng vào các bài toán thực tế), Câu 22"
%% kien_thuc: "Bài 8 (tọa độ vectơ, tích vectơ với một số), Bài 6 (tổng các lực cân bằng)"
%% trang_thai: chinh-thuc
%% ghi_chu: "Đề gốc dùng chung giả thiết cho Câu 22–23; đã chép vào từng câu và nói rõ mỗi lực có phương dọc theo chân máy"
\begin{ex}
Một chiếc máy quay được đặt trên một giá đỡ ba chân. Điểm đặt máy là $E(0;0;6)$. Các điểm tiếp xúc với mặt đất của ba chân lần lượt là
\[A_1(0;1;0),\quad A_2\left(-\dfrac{\sqrt3}2;-\dfrac12;0\right),\quad A_3\left(\dfrac{\sqrt3}2;-\dfrac12;0\right)\]
(đơn vị: mét). Trọng lực tác dụng lên máy quay là $\vec P$ có độ lớn $300$ N hướng thẳng đứng xuống dưới mặt đất $(Oxy)$. Gọi $\vec F_1$, $\vec F_2$, $\vec F_3$ là lực do ba chân máy tác dụng lên điểm $E$ để giữ hệ thống cân bằng (mỗi lực có phương dọc theo chân máy tương ứng).
Tọa độ của vectơ lực $\vec F_2$ là
\choice{\True $\vec F_2=\left(\dfrac{25}{\sqrt3};\dfrac{25}3;100\right)$}{$\vec F_2=\left(-\dfrac{25}{\sqrt3};\dfrac{25}3;100\right)$}{$\vec F_2=\left(\dfrac{25}{\sqrt3};-\dfrac{25}3;100\right)$}{$\vec F_2=\left(-\dfrac{25}{\sqrt3};-\dfrac{25}3;-100\right)$}
\loigiai{Do đối xứng, $\vec F_i=k\overrightarrow{A_iE}$ với cùng $k>0$; $\vec F_1+\vec F_2+\vec F_3=-\vec P=(0;0;300)$ nên $3\cdot6k=300$, $k=\dfrac{50}3$.

$\vec F_2=\dfrac{50}3\left(\dfrac{\sqrt3}2;\dfrac12;6\right)=\left(\dfrac{25}{\sqrt3};\dfrac{25}3;100\right)$.}
\end{ex}
